\chapter{The Misallocation of Intelligence}
\label{ch:misallocation}

The contemporary debate over centralized artificial intelligence infrastructure frequently rests on a premise that is seldom stated outright: that a scarce computational resource is used well to the extent that it is used continuously. On this premise, memory installed in personal computers is a stranded asset whenever it is not serving a request, and the natural remedy is to aggregate it, schedule it, and bring its utilization closer to that of a purpose-built facility. The argument was stated plainly by Sam Altman in an October 2026 podcast appearance, as reported by Bonifield~\cite{bonifield2026}. Asked about the ongoing shortage of DRAM, Altman said that from a global allocation perspective it is quite terrible for everyone to hold a share of the world's limited memory supply in a personal computer where it is used five percent of the time, and that maximum utilization in the cloud is the right aim in a supply-constrained environment. The argument has real force. Hardware that is manufactured and then left unused represents embodied energy and material that produced no service.

The question this chapter poses is whether maximizing the utilization of a scarce resource is equivalent to maximizing its social value. The two coincide only under conditions that rarely hold together: that the resource has a single use, that its idle time has no value of its own, that the cost of coordinating its use is negligible, and that the measure of utilization tracks the purpose for which the resource exists. Where any of these fails, a high utilization figure and a high social value can diverge, and the divergence is invisible to anyone who consults only the figure.

\section{Six Quantities That Are Not One}

The vocabulary of efficiency conceals several distinct quantities. Instantaneous utilization is the fraction of capacity engaged at a given moment. Lifetime utilization is the fraction of capacity engaged over the service life of the hardware, which can be high even when instantaneous utilization is usually low, provided the bursts are long enough. Economic efficiency compares the value of outputs to the value of all inputs, and depends on prices that are themselves shaped by institutional arrangements. Energy efficiency compares useful output to energy consumed, and says nothing about the energy embodied in manufacture. Resilience is the ability to continue to provide a service under disturbance, and it is typically purchased with redundancy that utilization metrics register as waste. Computational autonomy is the ability of an owner to run calculations without depending on a third party's network, terms, or continued existence.

These six quantities cannot be collapsed into a single objective without choosing weights, and the choice of weights is a value judgment rather than a technical result. A system tuned for instantaneous utilization will tend to sacrifice resilience and autonomy, because both require capacity held in reserve. A system tuned for autonomy will tend to show poor utilization, because the owner's machine must be able to serve the owner at any moment. The tension between them is structural. It is not a transitional inefficiency awaiting a better algorithm.

\section{Productive and Wasteful Idleness}

A computer that is idle for most of the day may still be doing something valuable: it is available. Availability has a price in any market that sells it. A standby generator, a fire engine, and a spare tire are all idle almost all of their lives, and no one describes their idleness as a failure of allocation. Their value lies in the option they confer, the ability to act immediately and locally when need arises.

This motivates a distinction between productive idleness and wasteful idleness. Idleness is productive when unused capacity functions as insurance against future demand, as protection of privacy or local control, or as headroom that keeps the primary function responsive. It is wasteful when no such function is served and the capacity could be directed elsewhere without cost to the owner. The difficulty is that the two cannot be told apart from a utilization trace alone. Both appear as a flat line near zero. Distinguishing them requires knowing what the idle capacity is for, which is information that resides with the owner and not with the scheduler.

\begin{definition}[Residual capability]
\label{def:residual}
Let a system $S$ serve a primary purpose $P$. A \emph{residual capability} of $S$ is a computational or physical action that $S$ can perform without degrading its performance of $P$ below a level specified independently of any use of the residual capability.
\end{definition}

The qualifier ``specified independently'' carries the weight of the definition. If the acceptable level of degradation is set by the party that wishes to use the residual capability, the definition becomes circular, and any capability whatsoever can be called residual by choosing a sufficiently tolerant standard. The same problem recurs in Chapter~\ref{ch:goodhart} as the general difficulty of metrics that are controlled by those they measure, and in Chapter~\ref{ch:idle} as the governance question of who decides what a system may spare.

\section{A Worked Comparison}

The inference in the utilization argument has three steps, and they are not equally secure. Higher utilization is taken to imply better allocation, and better allocation is taken to imply the preferred social arrangement. The first step requires empirical assumptions about workload, memory per unit of workload, and the cost of whatever infrastructure delivers the service. The second requires a normative argument that a utilization figure cannot supply, because it concerns who controls the resource and on what terms.

The first step can be made explicit. Let $W$ denote workload demand, $m$ the memory required per unit of workload, $u$ the utilization achieved, and $I$ the additional infrastructure needed to deliver the service, all expressed in comparable resource units. Writing subscript $c$ for centralized provision and $l$ for local provision, total resource consumption is approximately
\[
R_c=\frac{W_c\,m_c}{u_c}+I_c,\qquad
R_l=\frac{W_l\,m_l}{u_l}+I_l.
\]
The condition $u_c>u_l$ does not establish $R_c<R_l$. Centralization adds networking, redundancy, cooling, security, and administration to $I_c$. It can also raise $W_c$, since a service that is easier to reach tends to be used more, so that the savings from higher utilization are partly or wholly absorbed by increased demand. The comparison that matters is the total resource consumed to deliver an equivalent level of service, including reliability and latency, and a single utilization figure leaves that comparison open.

Two further observations bear on how the argument should be weighed. The reported figure of five percent is given without a stated methodology, and memory occupancy, memory bandwidth, and computational utilization are different quantities. And the scarcity to which the argument responds is partly the product of the same allocation. The report states that large memory manufacturers hold long-term agreements with cloud operators and with projects such as Stargate that have reserved much of the world's output for the next several years~\cite{bonifield2026}. A scarcity shaped by the purchases of the parties advocating centralization cannot simply be treated as an external constraint, and those parties also benefit commercially if consumers shift toward cloud inference. This does not make the engineering observation false. It means that a proposal to allocate scarce resources toward the proposer's business model should be assessed on the total-resource comparison, not on the proposer's chosen metric.

\section{Terms Used Throughout}

Four terms recur and are kept distinct, because each can hold while the next fails. Available capacity is what a device can physically do beyond its present load. It describes hardware and implies no permission. A residual capability, in the sense of Definition~\ref{def:residual}, is the narrower notion of an action the device can perform without degrading its primary function below a level specified independently of the use. Authorized capacity is the part of that which the owner permits and the safety envelope allows, as developed in Chapter~\ref{ch:idle}. Useful computation is computation whose results have independently verified value, as developed in Chapter~\ref{ch:norewards}. A contribution is admissible, in the sense of Definition~\ref{def:admissible}, only when it is residual, authorized, verified, and valuable. A resource can be physically idle without being authorized. An authorized contribution can still be economically disadvantageous. Completed computation can lack independently verified value. A statement in a later chapter that uses one of these terms should be read with its own meaning, and a result about one is not a result about the others.

\section{The Plan of the Argument}

The remaining chapters of Part~I develop the philosophical and institutional background to this problem. Chapter~\ref{ch:ellul} examines the claim that the pursuit of efficiency, once elevated to a supreme criterion, subordinates the purposes it was meant to serve. Chapter~\ref{ch:goodhart} supplies a formal account of how a measurable proxy can diverge from the objective it stands for under optimization pressure. Chapter~\ref{ch:bitcoin} considers proof-of-work mining as the clearest existing case of a reward structure that converts computational expenditure into an end.

Part~II then turns to an alternative organization of computation in which the unit of account is not utilization but the admissible use of capacity that already serves another purpose. The alternative may prove less efficient at first, and it may prove less efficient throughout. The monograph treats it as a research hypothesis whose case, if it can be made, rests on avoided capital, local authority, and learning over a long horizon. The later chapters test the hypothesis against findings that cut the other way. Centralized operators can learn faster in calendar time (Chapter~\ref{ch:horizon}). Distributed capacity is geographically correlated (Chapter~\ref{ch:geography}). Displacing heat pumps can raise peak electricity demand (Chapter~\ref{ch:geography}). Existing hardware is not free to operate (Chapter~\ref{ch:sunk}). And the alternative is only partly resistant to the failure mode that Part~I describes (Chapter~\ref{ch:norewards}).
